% !TeX program = lualatex
% ================================================================
% Polish Tier 3 Vocabulary Version of PieChartsAndProtractors
%
% Generated by the server using the following settings:
%
%   language            Polish (pol)
%   text direction      left to right
%   font                Noto Serif
%   fallback font       Noto Serif
%   built from          latex/worksheets/charts/pieChartsAndProtractors.tex
%   vocabulary key      latex/worksheets/charts/pieChartsAndProtractors/L2/pol/pieChartsAndProtractors_polishKey.tex
%   tier 3 vocabulary   green   RGB 0 128 0
%   English text        black   RGB 0 0 0
%   Polish text         purple  RGB 112 48 160
%   layout macros       version 3
%
% Please don't edit any information in this comment block - the server
% compares it against the current settings and rebuilds this file when
% they differ, so changes here would be overwritten. However, feel free
% to make updates to the LaTeX below if you want to edit it.
% ================================================================

\documentclass[a4paper,12pt]{article}
% ================================================================
% Copied from the English sheet this was translated from, so that
% anything it sets up for itself still works here
% ================================================================
\usepackage[a4paper, margin=1.8cm]{geometry}
\usepackage{tikz}
\usepackage{amsmath}
\usepackage{amssymb}
\usepackage{array}
\usepackage{mdframed}
\usepackage{enumitem}
\usepackage[T1]{fontenc}
\usepackage{textcomp}
\usetikzlibrary{calc}

% ---- Colours ----
\definecolor{slice1}{RGB}{70,130,180}
\definecolor{slice2}{RGB}{240,140,40}
\definecolor{slice3}{RGB}{70,170,70}
\definecolor{slice4}{RGB}{210,70,70}
\definecolor{slice5}{RGB}{150,90,190}
\definecolor{lightgrey}{RGB}{245,245,245}
\definecolor{boxblue}{RGB}{220,235,250}
\definecolor{boxgreen}{RGB}{220,245,220}

% -------------------------------------------------------
% \piechart{radius_number}{Label/Angle/Colour, ...}
%   radius_number is a plain number interpreted as cm.
%   Angles in degrees, clockwise from 12 o'clock; must sum to 360.
% -------------------------------------------------------
\newcommand{\piechart}[2]{%
  \begin{tikzpicture}
    \pgfmathsetmacro{\r}{#1}%
    \pgfmathsetmacro{\startA}{90}%
    \foreach \lbl/\ang/\col in {#2} {%
      \pgfmathsetmacro{\endA}{\startA - \ang}%
      \fill[\col!20] (0,0) -- (\startA:\r cm) arc(\startA:\endA:\r cm) -- cycle;%
      \draw[line width=2pt] (0,0) -- (\startA:\r cm);%
      \pgfmathsetmacro{\midA}{(\startA + \endA)/2}%
      \pgfmathsetmacro{\lR}{\r * 0.65}%
      \node[font=\small\bfseries, align=center] at (\midA:\lR cm) {\lbl};%
      \xdef\startA{\endA}%
    }%
    \draw[black, line width=1.8pt] (0,0) circle (\r cm);%
    % Draw all dividing lines again on top, including 12-o'clock
    \draw[black, line width=2pt] (0,0) -- (90:\r cm);%
    \fill[black] (0,0) circle (3pt);%
    \draw[black, line width=1.2pt] (0,\r cm) -- (0,{\r cm+0.25cm});%
  \end{tikzpicture}%
}

% Shortcut for answer boxes
\newcommand{\ansbox}{\underline{\hspace{1.6cm}}}
\newcommand{\sqfrac}{\dfrac{\square}{\square}}

\pagestyle{empty}

% Characters this language's font has not got are borrowed from another,
% rather than being silently dropped - see docs/translated-sheet-latex.md
\directlua{%
  if luaotfload and luaotfload.add_fallback then
    luaotfload.add_fallback("eallatinfallback",
      { "Noto Serif:mode=harf;" })
  else
    tex.error("this luaotfload is too old to register a fallback font")
  end
}
\usepackage[english]{babel}
\babelprovide[import]{polish}
\babelfont[polish]{rm}[Renderer=Harfbuzz, BoldFont={Noto Serif}, ItalicFont={Noto Serif}, BoldItalicFont={Noto Serif}, RawFeature={fallback=eallatinfallback}]{Noto Serif}
\babelfont[polish]{sf}[Renderer=Harfbuzz, BoldFont={Noto Serif}, ItalicFont={Noto Serif}, BoldItalicFont={Noto Serif}, RawFeature={fallback=eallatinfallback}]{Noto Serif}
\newcommand{\ealtext}[1]{\foreignlanguage{polish}{#1}}
\newcommand{\ealtextblock}[1]{{\raggedright\foreignlanguage{polish}{#1}\par}}
% ================================================================
% EAL layout helpers
% ================================================================
\usepackage{xcolor}

\definecolor{ealkeycolour}{RGB}{0,128,0}
\definecolor{ealenglish}{RGB}{0,0,0}
\definecolor{ealtwocolour}{RGB}{112,48,160}

\newsavebox{\ealboxen}
\newsavebox{\ealboxtr}
\newlength{\ealglosswd}
\newlength{\ealglossraise}
\setlength{\ealglossraise}{1.15em}

% a tier 3 word, in English and in translation alike
\newcommand{\ealkey}[1]{\textcolor{ealkeycolour}{#1}}
\newcommand{\ealkeytr}[1]{\textcolor{ealkeycolour}{\ealtext{#1}}}

% One translated word above one English word. Built from kernel
% primitives, never a tabular, which would break wherever the sheet
% loads array, booktabs or tabularx. The box is as wide as the wider of
% the two, so a long translation pushes its neighbours aside instead of
% overlapping them, and it claims its full height so a table row or a
% TikZ node makes room for it.
\newcommand{\ealgloss}[2]{%
  \sbox{\ealboxen}{\ealkey{#1}}%
  \sbox{\ealboxtr}{\small\textcolor{ealtwocolour}{\ealtext{#2}}}%
  \setlength{\ealglosswd}{\wd\ealboxen}%
  \ifdim\wd\ealboxtr>\ealglosswd\setlength{\ealglosswd}{\wd\ealboxtr}\fi
  \leavevmode
  \makebox[\ealglosswd][c]{%
    \makebox[0pt][c]{\usebox{\ealboxen}}%
    \makebox[0pt][c]{\raisebox{\ealglossraise}{\usebox{\ealboxtr}}}%
  }%
}

% opens up line spacing around a block of glosses, so the raised
% translations sit clear of the line above
\newenvironment{ealglossed}{\par\linespread{2.1}\selectfont\ignorespaces}{\par}

% a whole translated sentence above its English counterpart. block level,
% so it does not belong inside a tabular cell
\newcommand{\ealpara}[2]{%
  \par\addvspace{0.5\baselineskip}%
  {\color{ealtwocolour}\ealtextblock{#1}}%
  \penalty200
  {\color{ealenglish}#2\par}%
  \addvspace{0.5\baselineskip}%
}

\begin{document}

% ===========================
%  TITLE
% ===========================
\begin{ealglossed}
\begin{center}
  {\LARGE\bfseries Reading \ealgloss{Pie Charts}{diagramy kołowe} with a \ealgloss{Protractor}{kątomierz}}\\[4pt]
  \vspace{2em}
  {Name:\;\underline{\hspace{6cm}} \quad
   Date:\;\underline{\hspace{2.8cm}}}
\end{center}
\end{ealglossed}
\vspace{4pt}

\begin{ealglossed}
\begin{mdframed}[backgroundcolor=boxgreen, linecolor=slice3,
                 linewidth=1.5pt, roundcorner=6pt, innertopmargin=6pt]
\textbf{Key \ealgloss{formulae}{wzory}:}\quad
$\displaystyle\text{sector fraction} = \frac{\text{sector angle}}{360}$
\qquad\qquad
$\displaystyle\text{amount} = \text{fraction} \times \text{total}$
\end{mdframed}
\end{ealglossed}
\vspace{8pt}

% ===========================
%  PIE CHART 1
% ===========================
\begin{ealglossed}
\begin{center}
  {\large\bfseries \ealgloss{Pie Chart}{diagram kołowy} 1 \textendash{} Favourite School Subjects (30 students)}
\end{center}
\end{ealglossed}

\begin{ealglossed}
\noindent
\begin{minipage}[c]{0.6\textwidth}
  \centering
  \piechart{5.8}{Maths/120/slice1,Science/90/slice2,English/72/slice3,Art/48/slice4,PE/30/slice5}
\end{minipage}%
\hfill
\begin{minipage}[c]{0.3\textwidth}
  \setlength{\tabcolsep}{4pt}
  \renewcommand{\arraystretch}{2.5}
  \begin{tabular}{|l|c|c|c|}
  \hline
  \textbf{Subject} & \textbf{\ealgloss{Angle}{kąt}} & \textbf{\ealgloss{Fraction}{ułamek}} & \textbf{No.} \\
  \hline
  Maths    & & & \\
  \hline
  Science  & & & \\
  \hline
  English  & & & \\
  \hline
  Art      & & & \\
  \hline
  PE       & & & \\
  \hline
  \textbf{\ealgloss{Total}{suma}} & \textbf{360\textdegree} & \textbf{1} & \textbf{30} \\
  \hline
  \end{tabular}
\end{minipage}
\end{ealglossed}

\begin{ealglossed}
\vspace{6pt}
\noindent\textbf{Questions -- \ealgloss{Pie Chart}{diagram kołowy} 1:}
\begin{enumerate}[leftmargin=1.6em, itemsep=4pt]
  \item Use your \ealgloss{protractor}{kątomierz} to measure each \ealgloss{sector}{wycinek koła} and fill in the \textbf{\ealgloss{Angle}{kąt}} column above.
  \item Do your \ealgloss{angles}{kąty} add up to 360\textdegree? Write your \ealgloss{total}{suma} here: \ansbox\textdegree
  \item Write the \ealgloss{fraction}{ułamek} of students who chose \textbf{Maths} in its \ealgloss{simplest form}{najprostsza postać}:
        \quad $\dfrac{\framebox{\phantom{120}}}{360} = \frac{\phantom{ABC}}{\phantom{ABC}}$
  \item How many students chose \textbf{Science}?
        \quad $\dfrac{\framebox{\phantom{120}}}{360} \times 30 = {}$\ansbox{} students
  \item Which subject was chosen by exactly $\dfrac{1}{5}$ of the class? \quad \ansbox
\end{enumerate}
\end{ealglossed}

\newpage

% ===========================
%  PIE CHART 2
% ===========================
\begin{ealglossed}
\begin{center}
  {\large\bfseries \ealgloss{Pie Chart}{diagram kołowy} 2 \textendash{} How Students Travel to School (60 students)}
\end{center}
\end{ealglossed}

\begin{ealglossed}
\noindent
\begin{minipage}[c]{0.6\textwidth}
  \centering
  \piechart{5.8}{Walk/90/slice3,Bus/120/slice1,Car/60/slice2,Cycle/45/slice4,Other/45/slice5}
\end{minipage}%
\hfill
\begin{minipage}[c]{0.3\textwidth}
  \setlength{\tabcolsep}{4pt}
  \renewcommand{\arraystretch}{2.5}
  \begin{tabular}{|l|c|c|c|}
  \hline
  \textbf{Method} & \textbf{\ealgloss{Angle}{kąt}} & \textbf{\ealgloss{Fraction}{ułamek}} & \textbf{No.} \\
  \hline
  Walk   & & & \\
  \hline
  Bus    & & & \\
  \hline
  Car    & & & \\
  \hline
  Cycle  & & & \\
  \hline
  Other  & & & \\
  \hline
  \textbf{\ealgloss{Total}{suma}} & \textbf{360\textdegree} & \textbf{1} & \textbf{60} \\
  \hline
  \end{tabular}
\end{minipage}
\end{ealglossed}

\begin{ealglossed}
\vspace{6pt}
\noindent\textbf{Questions -- \ealgloss{Pie Chart}{diagram kołowy} 2:}
\begin{enumerate}[leftmargin=1.6em, itemsep=4pt]
  \item Measure all \ealgloss{sectors}{wycinki koła} with your \ealgloss{protractor}{kątomierz}. Record the \ealgloss{angles}{kąty} in the table above.
  \item How many students travel by \textbf{Bus}?
        \quad $\dfrac{\framebox{\phantom{120}}}{360} \times 60 = {}$\ansbox{} students
  \item Write the \ealgloss{fraction}{ułamek} of students who \textbf{walk} in its \ealgloss{simplest form}{najprostsza postać}:
        \quad $\dfrac{\framebox{\phantom{120}}}{360} = \frac{\phantom{ABC}}{\phantom{ABC}}$
  \item Two \ealgloss{sectors}{wycinki koła} have the \textbf{same} \ealgloss{angle}{kąt}. Which two travel methods are they? \quad \ansbox
  \item How many \textbf{more} students travel by Bus than by Car? \quad \ansbox{} students
\end{enumerate}
\end{ealglossed}

\newpage

% ===========================
%  PIE CHART 3
% ===========================
\begin{ealglossed}
\begin{center}
  {\large\bfseries \ealgloss{Pie Chart}{diagram kołowy} 3 \textendash{} Ice Cream Flavours Sold in One Day (180 sold)}
\end{center}
\end{ealglossed}

\begin{ealglossed}
\noindent
\begin{minipage}[c]{0.6\textwidth}
  \centering
  \piechart{5.8}{Vanilla/80/slice2,Choc/100/slice1,Straw/60/slice4,Mint/40/slice3,Other/80/slice5}
\end{minipage}%
\hfill
\begin{minipage}[c]{0.3\textwidth}
  \setlength{\tabcolsep}{4pt}
  \renewcommand{\arraystretch}{2.5}
  \begin{tabular}{|l|c|c|c|}
  \hline
  \textbf{Flavour} & \textbf{\ealgloss{Angle}{kąt}} & \textbf{\ealgloss{Fraction}{ułamek}} & \textbf{No.} \\
  \hline
  Vanilla    & & & \\
  \hline
  Chocolate  & & & \\
  \hline
  Strawberry & & & \\
  \hline
  Mint       & & & \\
  \hline
  Other      & & & \\
  \hline
  \textbf{\ealgloss{Total}{suma}} & \textbf{360\textdegree} & \textbf{1} & \textbf{180} \\
  \hline
  \end{tabular}
\end{minipage}
\end{ealglossed}

\begin{ealglossed}
\vspace{6pt}
\noindent\textbf{Questions -- \ealgloss{Pie Chart}{diagram kołowy} 3:}
\begin{enumerate}[leftmargin=1.6em, itemsep=4pt]
  \item Measure all \ealgloss{sectors}{wycinki koła}. Check your \ealgloss{angles}{kąty} \ealgloss{sum}{suma} to \textbf{360\textdegree} before moving on.
  \item How many \textbf{Chocolate} ice creams were sold?
        \quad $\dfrac{\framebox{\phantom{120}}}{360} \times 180 = {}$\ansbox{} ice creams
  \item How many \textbf{Mint} ice creams were sold? \quad \ansbox{} ice creams
  \item \textbf{Vanilla} and \textbf{Other} have the same \ealgloss{angle}{kąt}. How many does each represent?
        \quad \ansbox{} ice creams each
  \item Which flavour is exactly $\tfrac{1}{6}$ of all sales? Show your working below.

        \vspace{1.5cm}
        Answer: \ansbox
\end{enumerate}
\end{ealglossed}

\newpage

% ===========================
%  PIE CHART 4
% ===========================
\begin{ealglossed}
\begin{center}
  {\large\bfseries \ealgloss{Pie Chart}{diagram kołowy} 4 \textendash{} One Evening's Activities (4 hours = 240 minutes)}
\end{center}
\end{ealglossed}

\begin{ealglossed}
\noindent
\begin{minipage}[c]{0.6\textwidth}
  \centering
  \piechart{5.8}{HW/90/slice4,TV/120/slice2,Gaming/60/slice1,Reading/30/slice3,Other/60/slice5}
\end{minipage}%
\hfill
\begin{minipage}[c]{0.31\textwidth}
  \setlength{\tabcolsep}{4pt}
  \renewcommand{\arraystretch}{2.5}
  \begin{tabular}{|l|c|c|c|}
  \hline
  \textbf{Activity} & \textbf{\ealgloss{Angle}{kąt}} & \textbf{\ealgloss{Fraction}{ułamek}} & \textbf{Mins} \\
  \hline
  Homework & & & \\
  \hline
  TV       & & & \\
  \hline
  Gaming   & & & \\
  \hline
  Reading  & & & \\
  \hline
  Other    & & & \\
  \hline
  \textbf{\ealgloss{Total}{suma}} & \textbf{360\textdegree} & \textbf{1} & \textbf{240} \\
  \hline
  \end{tabular}
\end{minipage}
\end{ealglossed}

\begin{ealglossed}
\vspace{6pt}
\noindent\textbf{Questions -- \ealgloss{Pie Chart}{diagram kołowy} 4:}
\begin{enumerate}[leftmargin=1.6em, itemsep=4pt]
  \item Measure all \ealgloss{angles}{kąty}. Complete the table. \textit{(Hint: 4 hours = 240 minutes.)}
  \item How many minutes were spent watching \textbf{TV}?
        \quad $\dfrac{\framebox{\phantom{120}}}{360} \times 240 = {}$\ansbox{} mins
  \item How many minutes were spent on \textbf{Homework}? \quad \ansbox{} mins
  \item Write the \ealgloss{fraction}{ułamek} of the evening spent \textbf{Gaming} in its \ealgloss{simplest form}{najprostsza postać}:
        \quad $\dfrac{\framebox{\phantom{120}}}{360} = \frac{\phantom{ABC}}{\phantom{ABC}}$
  \item \textbf{Challenge:} Another student spends the same \ealgloss{fraction}{ułamek} of their evening on Homework,
        but has \textbf{5 hours} (300 minutes) free. How many minutes do they spend on homework?

        \vspace{0.8cm}
        Answer: \ansbox{} mins
\end{enumerate}
\end{ealglossed}

\newpage{}
% ===========================
%  ANSWER KEY
% ===========================
\begin{mdframed}[backgroundcolor=lightgrey, linecolor=gray,
                 linewidth=1pt, roundcorner=4pt, innertopmargin=5pt]
\textbf{\small Answers (Teacher Copy):}\\[3pt]
{\small
\textbf{Chart 1:}
Maths 120\textdegree\;$=\tfrac{1}{3}=10$;\;
Science 90\textdegree\;$=\tfrac{1}{4}=9$;\;
English 72\textdegree\;$=\tfrac{1}{5}=6$;\;
Art 48\textdegree\;$=\tfrac{2}{15}=4$;\;
PE 30\textdegree\;$=\tfrac{1}{12}=3$.
\;Q3: $\tfrac{1}{3}$.\;Q4: 9 students.\;Q5: English.\\[3pt]
\textbf{Chart 2:}
Walk 90\textdegree\;$=15$;\;
Bus 120\textdegree\;$=20$;\;
Car 60\textdegree\;$=10$;\;
Cycle 45\textdegree\;$\approx8$;\;
Other 45\textdegree\;$\approx8$.
\;Q2: 20.\;Q3: $\tfrac{1}{4}$.\;Q4: Cycle \& Other.\;Q5: 10 more.\\[3pt]
\textbf{Chart 3:}
Vanilla 80\textdegree\;$=40$;\;
Choc 100\textdegree\;$=50$;\;
Straw 60\textdegree\;$=30$;\;
Mint 40\textdegree\;$=20$;\;
Other 80\textdegree\;$=40$.
\;Q2: 50.\;Q3: 20.\;Q4: 40 each.\;Q5: Strawberry ($60\textdegree=\tfrac{1}{6}$).\\[3pt]
\textbf{Chart 4:}
HW 90\textdegree\;$=60$\,min;\;
TV 120\textdegree\;$=80$\,min;\;
Gaming 60\textdegree\;$=40$\,min;\;
Reading 30\textdegree\;$=20$\,min;\;
Other 60\textdegree\;$=40$\,min.
\;Q2: 80\,min.\;Q3: 60\,min.\;Q4: $\tfrac{1}{6}$.\;Challenge: $\tfrac{1}{4}\times300=75$\,min.
}
\end{mdframed}

\end{document}